\section{Оптимальное управление}
\label{sec:optimal-control}

\subsection{Постановка задачи оптимального управления}

\emph{Оптимальное управление} изучает выбор управляющего воздействия, при
котором динамическая система переводится в требуемое состояние и одновременно
достигается наилучшее значение заданного критерия качества. Управляемый объект
обычно описывается системой обыкновенных дифференциальных уравнений
\[
    \dot X=f(X,U,t),\qquad X(t_0)=X_0,
\]
где $X(t)\in\mathbb R^n$ --- вектор состояния, а
$U(t)\in\mathbb R^m$ --- вектор управляющих воздействий.

Для управления могут быть заданы ограничения, например
\[
    U(t)\in\mathcal U,
    \qquad |u_i(t)|\leq U_i^{\max}.
\]
Управление, удовлетворяющее ограничениям и обеспечивающее существование
соответствующей траектории, называется \emph{допустимым}.

С управлением $U(t)$ однозначно связывается траектория $X(t)$. Пара
$(X^*(t),U^*(t))$, на которой критерий качества принимает наилучшее значение,
называется оптимальным процессом.

\subsection{Функционал качества}

В простейшем случае качество процесса задаётся интегральным функционалом
\[
    J(X,U)=\int_{t_0}^{t_f} f^0(X(t),U(t),t)\,dt
    \longrightarrow \min.
\]
Более общий функционал Больца включает также терминальную часть:
\[
    J(X,U)=\Phi(X(t_f))+
    \int_{t_0}^{t_f} f^0(X(t),U(t),t)\,dt.
\]
Функционал может характеризовать, например, затраты энергии на управление,
отклонение состояния от требуемого или время переходного процесса.

Типичный квадратичный критерий имеет вид
\[
    J=\int_{t_0}^{t_f}
    \left(X^TQX+U^TRU\right)dt,
\]
где $Q\geq0$, $R>0$ --- весовые матрицы. Первый член штрафует отклонение
состояния, второй --- чрезмерно большие управляющие воздействия.

Различают \emph{программное} управление $U=U(t)$, заранее заданное как функция
времени, и управление с обратной связью $U=U(X,t)$, которое корректируется по
текущему состоянию системы.

\subsection{Гамильтониан и необходимые условия оптимальности}

Для задачи
\[
    \dot X=f(X,U,t),\qquad
    J=\int_{t_0}^{t_f}f^0(X,U,t)\,dt\to\min
\]
введём сопряжённый вектор $\lambda(t)$. В общей формулировке принципа
максимума вводится также скалярный множитель $\lambda_0\geq0$ и гамильтониан
\[
    H(X,U,\lambda,\lambda_0,t)
    =-\lambda_0 f^0(X,U,t)+\lambda^Tf(X,U,t),
\]
причём пара множителей $(\lambda_0,\lambda(t))$ не должна быть тождественно
нулевой. В \emph{нормальном случае} $\lambda_0>0$ его можно нормировать:
$\lambda_0=1$. Далее используется именно эта нормировка,
\[
    H(X,U,\lambda,t)
    =-f^0(X,U,t)+\lambda^Tf(X,U,t).
\]
Тогда необходимые условия оптимальности принимают канонический вид
\[
    \dot X=\left(\frac{\partial H}{\partial\lambda}\right)^T,
    \qquad
    \dot\lambda=-\left(\frac{\partial H}{\partial X}\right)^T.
\]
Если оптимальное управление является внутренней точкой допустимого множества и
гамильтониан гладок по $U$, выполняется условие стационарности
\[
    \frac{\partial H}{\partial U}=0.
\]
Оно позволяет выразить управление через состояние и сопряжённые переменные:
\[
    U=U^*(X,\lambda,t).
\]
Таким образом, вместо исходной задачи оптимизации получается система уравнений
для $X$ и $\lambda$ с условиями на левом и правом концах, то есть, как правило,
двухточечная краевая задача.

Если конечное состояние фиксировано, задаётся условие $X(t_f)=X_f$. Если
конечное состояние свободно, возникают условия трансверсальности. Например,
для функционала с терминальной частью $\Phi(X(t_f))$ при свободном $X(t_f)$
сопряжённый вектор на правом конце определяется производной терминального
слагаемого с учётом выбранного соглашения знаков.

\subsection{Принцип максимума Понтрягина}

При наличии ограничений на управление условие
$\partial H/\partial U=0$ применять непосредственно нельзя: максимум может
достигаться на границе множества допустимых управлений. Основным необходимым
условием становится \emph{принцип максимума Понтрягина}.

Если $(X^*(t),U^*(t))$ --- оптимальный процесс, то существуют нетривиальные
множители Понтрягина; в рассматриваемом нормальном случае после нормировки
$\lambda_0=1$ сопряжённый вектор $\lambda(t)$ удовлетворяет каноническим
уравнениям и почти для всех $t$
\[
    H(X^*(t),U^*(t),\lambda(t),t)
    =\max_{U\in\mathcal U}
      H(X^*(t),U,\lambda(t),t).
\]
Следовательно, оптимальное управление в каждый момент выбирается из допустимого
множества так, чтобы максимизировать гамильтониан.

Важно, что принцип максимума даёт в общем случае \emph{необходимые}, а не
достаточные условия оптимальности. Поэтому найденную экстремаль необходимо
сопоставлять с граничными условиями и исходным критерием качества.

Если система и критерий не зависят от времени явно и конечное время фиксировано,
то вдоль оптимальной траектории гамильтониан постоянен. При свободном конечном
времени в стандартной постановке возникает дополнительное условие на значение
гамильтониана в конечный момент.

\subsection{Линейно-квадратичная задача}

Один из важнейших специальных случаев --- линейная система
\[
    \dot X=AX+BU
\]
с квадратичным функционалом
\[
    J=\frac12 X^T(t_f)P_fX(t_f)+
    \int_{t_0}^{t_f}\left(X^TQX+\rho U^TRU\right)dt,
    \qquad \rho>0.
\]
Гамильтониан можно записать как
\[
    H=-X^TQX-\rho U^TRU+\lambda^T(AX+BU).
\]
Из условия стационарности следует
\[
    U^*(t)=\frac{1}{2\rho}R^{-1}B^T\lambda(t),
\]
а сопряжённое уравнение имеет вид
\[
    \dot\lambda=2QX-A^T\lambda.
\]
Тем самым задача снова сводится к совместному решению уравнений состояния и
сопряжённой системы. В задаче оптимальной стабилизации с обратной связью решение
приводит к линейному закону $U=-K(t)X$, а матрица усиления определяется через
решение матричного уравнения Риккати.

\subsection{Что важно сказать на экзамене}

Нужно начать с модели $\dot X=f(X,U,t)$ и функционала качества $J\to\min$,
объяснить, что ищется допустимое управление $U^*(t)$ и соответствующая
оптимальная траектория. Затем следует отметить общую форму $H=-\lambda_0f^0+\lambda^Tf$ и
нетривиальность пары множителей, после чего в нормальном случае положить
$\lambda_0=1$, записать уравнения состояния и сопряжённой системы и
сформулировать принцип максимума Понтрягина
\[
    H(X^*,U^*,\lambda,t)=\max_{U\in\mathcal U}H(X^*,U,\lambda,t).
\]
В качестве стандартного примера достаточно назвать линейно-квадратичную задачу,
для которой оптимальное управление определяется из условия стационарности, а
при построении регулятора возникает уравнение Риккати.

\newpage
